\section{Related Work}
\label{sec:related}

\paragraph{MCP security.}
The MCP specification~\citep{mcp2024} defines the transport, the tool invocation
protocol, and OAuth-based authorization of a client to a server, but leaves per-call policy
to implementors. Concurrent work has identified
prompt-injection and tool-poisoning risks in MCP deployments~\citep{mcpsecurity2025}, where
malicious tool descriptions or responses cause an agent to take unintended actions. cMCP
addresses tool-poisoning through catalog hash pinning (Property~\ref{prop:rug-pull}) but
does not claim to prevent all prompt-injection attacks; the inspection pipeline detects
known patterns but cannot make semantic judgments about agent intent.

\paragraph{TEE-based policy enforcement.}
Using TEEs to enforce policies at runtime has been studied in the context of database
systems~\citep{teedb2017}, open enclave frameworks~\citep{keystone2020}, and network
functions~\citep{teefw2019}. The common pattern is to measure a policy into the TEE at
startup and use remote attestation to let a verifier confirm which policy was in effect.
cMCP applies this pattern to the MCP boundary, adding the dimension of session-level state
(the sensitivity state machine) that persists across multiple tool calls within a session,
rather than evaluating each request independently.

\paragraph{Agent governance frameworks.}
The Agent Governance Toolkit (AGT)~\citep{agt2026} provides the policy engine, identity
layer (SPIFFE/SVID), and Cedar policy bundle format that cMCP uses for enforcement.
cMCP acts as an enforcement point for AGT policies at the MCP boundary. Other agent
governance frameworks~\citep{guardrails2023,nemo2023} focus on output filtering and input
validation rather than tool-call interception, and do not provide hardware-backed
attestation of enforcement.

\paragraph{Attestation standards.}
RATS (Remote ATtestation procedureS)~\citep{rats2021} defines the roles and information
flows for remote attestation: Attester (the device producing evidence), Verifier (who
appraises it), and Relying Party (who makes trust decisions). TRACE~\citep{trace2025}
implements the RATS model for AI agent deployments. The TRACE Claim produced by cMCP is
an Attestation Result in RATS terminology, combining hardware evidence (the TEE
measurement report) with application-level evidence (the session call summary and audit
chain). EAT (Entity Attestation Token)~\citep{eat2023} provides the underlying token
format that TRACE extends.

\paragraph{Cryptographic audit logs.}
Hash-chained audit logs are well-studied in the context of certificate
transparency~\citep{ct2013} and secure logging~\citep{seclog2000}. The cMCP audit chain
is a linear hash chain in the style of Schneier and Kelsey: each entry commits to the hash
of its predecessor, so any deletion or modification of an entry changes the chain tip,
which is embedded in the TRACE Claim. Certificate transparency instead uses a Merkle tree,
which gives logarithmic inclusion and consistency proofs that a linear chain does not.
cMCP chains are also per-session rather than global.

\paragraph{Data provenance.}
Provenance tracking in distributed systems has a long history~\citep{prov2013}; PROV-DM
provides a data model for expressing provenance relationships. Temporal adjacency, as used
in cMCP, is a weak form of provenance: it asserts that call $B$ \emph{could} be
PHI-adjacent because it followed a PHI call, but does not assert that $B$ \emph{actually}
carries PHI-derived data. This is analogous to conservative call graph approximations in
static analysis~\citep{reps1995}, which accept false edges to achieve soundness. The
provenance disclaimer in every cMCP call graph summary (Property~\ref{prop:temporal}(c))
makes this distinction explicit.

\paragraph{Confidential computing for AI.}
Prior work applies confidential computing to machine learning~\citep{confai2024}, and
cryptographic protocols such as homomorphic encryption protect transformer inference
without trusted hardware~\citep{teeinference2023}. cMCP's use of TEEs is orthogonal: it
protects the governance layer (policy enforcement, audit) rather than the model itself.
The two approaches are composable: a model running in a TEE and a cMCP gateway both
running in TEEs can exchange mutual attestation reports, establishing end-to-end hardware
trust across the full inference-to-tool-call path.
